\section{Certified radial and spectral consequences}
\label{sec:consequences}

Throughout this section, $F_\star:=F_*$ denotes the exact profile in the
certified Newton neighborhood, not the rational center spline, and
$I_{\rm rot}:=\mathcal I[F_*]$ is the dimensionless inertia defined in
\cref{thm:main}.  The two consequences below import the same authenticated
wall-slope and inertia enclosures,
\begin{align}
 F_\star'(4)&\in[-0.094644972770043,-0.087466485763471],\notag\\
 I_{\rm rot}&\in[21.149277305080,48.920186135018].
 \label{eq:consequence-profile-inputs}
\end{align}
They are therefore statements about the validated nonlinear solution.  The
combined profile-and-tail archive has SHA-256 digest
\begin{center}
\small\nolinkurl{1d5fe53786cc280006d7b1092d360556d4d8d8684e5ae3356ce8cd6d084e72a9}.
\end{center}
Decimal endpoints in this section are outward-rounded displays of rational
intervals; the archived fractions, rather than these displays, are used in the
proof comparisons.

\subsection{Fixed-wall spherical radial gap}
\label{sec:fixed-wall-gap}

Let $\tau=e f_\pi t$.  Thus a mode with dimensionless frequency
$\widehat\omega$ has Killing-time frequency
$\omega_K=e f_\pi\widehat\omega$.  Set
\begin{equation}
 u(F,x)=x^2+8\sin^2F,
 \qquad N(x)=1-\lambda x^2.
\end{equation}
Linearization of the spherically symmetric Skyrme action about $F_\star$ gives
the generalized Sturm--Liouville problem
\begin{equation}
 \ell_\star\eta
   =\widehat\omega^{2}\mathcal W_\star\eta,
 \qquad
 \ell_\star=-\frac{\dd}{\dd x}
       \left(P_\star\frac{\dd}{\dd x}\right)+Q_\star,
 \label{eq:radial-sturm-liouville}
\end{equation}
where
\begin{equation}
 P_\star=N u(F_\star,x),
 \qquad
 \mathcal W_\star=\frac{u(F_\star,x)}{N}.
 \label{eq:radial-weights}
\end{equation}
Here $Q_\star$ is precisely the second-variation potential of the static
functional used in the validated Jacobi problem.  To specify the singular
endpoint, let
\begin{equation}
 \mathscr H=L^2\bigl((0,4),\mathcal W_\star(x)\dd x\bigr)
\end{equation}
and initially define
\begin{equation}
 \mathfrak q_\star[\eta]
 =\int_0^4\left(P_\star|\eta'|^2+Q_\star|\eta|^2\right)\dd x
 \label{eq:radial-quadratic-form}
\end{equation}
on $C_c^\infty(0,4)$.  Its closure is the fixed-wall Friedrichs form.  The
associated self-adjoint operator in $\mathscr H$ is
\begin{equation}
 \mathcal A_\star=\mathcal W_\star^{-1}\ell_\star,
 \qquad \mathcal A_\star\eta=\widehat\omega^2\eta.
 \label{eq:radial-self-adjoint-operator}
\end{equation}
The regular-singular classification below shows that eigenmodes
select the physical branch $\eta(x)=O(x)$ at the origin; the closure enforces
the Dirichlet trace at $x=4$.

\begin{proposition}[Authenticated fixed-wall radial stability]
\label{prop:fixed-wall-gap}
For $\mu^2=1$, $\lambda=1/400$, and $x_{\rm w}=4$, the exact profile satisfies
\begin{equation}
 \mathfrak q_\star[\eta]
 \geq \norm{\eta}_{L^2(0,4)}^2
 \label{eq:static-form-gap}
\end{equation}
on the closed Friedrichs form domain.  The operator $\mathcal A_\star$ has compact
resolvent in $\mathscr H$, and its lowest dimensionless spherical profile
frequency obeys
\begin{equation}
 \widehat\omega_{\rm rad}^2\geq\frac1{25},
 \qquad
 \widehat\omega_{\rm rad}\geq\frac15,
 \qquad
 \omega_{K}\geq\frac{e f_\pi}{5}.
 \label{eq:fixed-wall-frequency-gap}
\end{equation}
\end{proposition}

\begin{proof}[Proof sketch]
Use the positive rational Barta witness
\begin{equation}
 v(x)=\frac{8}{(x-33/16)^2+4}.
\end{equation}
On $[1/16,4]$, adaptive interval evaluation through every exact-solution jet
in the Newton tube gives
\begin{equation}
 \inf_{[1/16,4]}\frac{\ell_\star v}{v}
 >1.0386099769.
 \label{eq:positive-radius-barta}
\end{equation}
The origin cannot be obtained by simply extending those boxes through the
regular singular point.  Instead, substituting
$\pi-F_\star=xu(x^2)$ and using the exact profile equation removes all negative
powers of $x$ and gives
\begin{equation}
 \inf_{[0,1/16]}\frac{\ell_\star v}{v}
 >36.8298881657.
 \label{eq:origin-barta}
\end{equation}
The ground-state transform with $v$ proves
\cref{eq:static-form-gap} first on $C_c^\infty(0,4)$ and then, by closure, on
the Friedrichs form domain.  The regular-origin expansion gives, with
$a_0=1+8b^2$,
\begin{equation}
 P_\star=a_0x^2+O(x^4),\qquad
 \mathcal W_\star=a_0x^2+O(x^4),\qquad
 Q_\star=2a_0+O(x^2).
 \label{eq:radial-origin-coefficients}
\end{equation}
Thus the indicial roots of the eigenvalue equation are $1$ and $-2$.
The $x^{-2}$ branch is not in $\mathscr H$, so the origin is limit point and
eigenmodes obey $\eta=O(x)$.  At $x=4$, where $P_\star>0$, the
Friedrichs closure gives the Dirichlet trace.

For completeness, the form embedding into $\mathscr H$ is compact.  On every
$[\delta,4]$, the ground-state representation controls an $H^1$ norm and
Rellich compactness applies.  Near the origin,
$\mathcal W_\star(x)\leq Cx^2$ and
\cref{eq:static-form-gap} give uniformly on form-bounded sets
\begin{equation}
 \int_0^\delta\mathcal W_\star|\eta|^2\dd x
 \leq C\delta^2\mathfrak q_\star[\eta].
\end{equation}
The vanishing tail and the interior compactness prove compact resolvent.

Finally, $N(4)=24/25$ and $0\leq\sin^2F_\star\leq1$, whence
\begin{equation}
 0\leq\mathcal W_\star(x)
 \leq\frac{x_{\rm w}^2+8}{N(x_{\rm w})}=25.
\end{equation}
The min--max principle for $\mathcal A_\star$, together with the generalized Rayleigh
quotient, converts the unit static form gap into
\cref{eq:fixed-wall-frequency-gap}.  The authenticated calculation used 109
positive-radius leaves, with maximum binary refinement depth five.
\end{proof}

The exact radial-gap archive has SHA-256 digest
\begin{center}
\small\nolinkurl{695310609d070f6dba2c678982ffc87472ed7c341813137fb90ad9dd9e1e6096}.
\end{center}
It verifies the identity of the upstream spectral/profile archive before
reconstructing the exact-solution tube.  In particular, it authenticates the
canonical spectral archive and sharp tube snapshot with respective digests
\begin{center}
\small
\nolinkurl{1d5fe53786cc280006d7b1092d360556d4d8d8684e5ae3356ce8cd6d084e72a9}\\
\nolinkurl{1781a2ff357f3b165d23e290eb403552d77d099a72b9e90920735e6a80b30431}.
\end{center}

\begin{remark}[Scope of the stability statement]
Proposition~\ref{prop:fixed-wall-gap} concerns only the $l=0$ profile
fluctuation with an ideal, nondynamical Dirichlet wall.  It is not a theorem
about a moving membrane, nonspherical or grand-spin sectors, the rotational
collective coordinate, a holding apparatus, or metric perturbations.  In
particular, it is a radial spectral gap on the declared fixed-wall form domain,
not full dynamical stability of a supported or self-gravitating Skyrmion.
\end{remark}

\subsection{Authenticated endpoint jets}
\label{sec:endpoint-jets}

Introduce the optical radius
\begin{equation}
 x(y)=\frac{\tanh y}{\sqrt\lambda},
 \qquad
 Y=\operatorname{arctanh}(\sqrt\lambda x_{\rm w}),
 \qquad
 N_{\rm w}=1-\lambda x_{\rm w}^2.
 \label{eq:optical-coordinate}
\end{equation}
Write the dimensionless inertia as
\begin{equation}
 I_{\rm rot}=\int_0^4\rho_I(x)\dd x,
 \qquad
 \rho_I(x)=\frac{2\pi}{3}x^2\sin^2F_\star
 \left[\frac1N+4(F_\star')^2
       +\frac{4\sin^2F_\star}{Nx^2}\right].
 \label{eq:inertia-density}
\end{equation}
Define the two optical weights
\begin{align}
 W(y)&=\frac{\rho_I(x(y))}{x(y)^2}\frac{\dd x}{\dd y}
 \notag\\
 &=\frac{2\pi\sin^2F_\star}{3\sqrt\lambda}
 \left[1+4\lambda\cosh^2y\,F_{\star,y}^2
       +4\lambda\coth^2y\,\sin^2F_\star\right],
 \qquad A(y)=W(y)\coth y.
 \label{eq:optical-weights}
\end{align}

At the regular origin,
\begin{align}
 F_\star(x)&=\pi-bx+cx^3+O(x^5),
 \notag\\
 c&=\frac{b\left[\mu^2-4\lambda+
       (\frac43-24\lambda)b^2+\frac83b^4\right]}
       {10(1+8b^2)},
 \label{eq:origin-profile-jet}\\
 W(y)&=w_0y^2+O(y^4),
 &A(y)&=w_0y+O(y^3),
 &w_0&=\frac{2\pi b^2(1+8b^2)}{3\lambda^{3/2}}.
 \label{eq:origin-optical-jets}
\end{align}
At the wall, let $\sigma=\abs{F_\star'(x_{\rm w})}$ and
$\delta=Y-y$.  Then
\begin{align}
 W(y)&=w_Y\delta^2+O(\delta^3),
 &A(y)&=\coth(Y)w_Y\delta^2+O(\delta^3),
 \label{eq:wall-optical-jets}\\
 w_Y&=\frac{2\pi\sigma^2N_{\rm w}^2
        (1+4N_{\rm w}\sigma^2)}{3\lambda^{3/2}},
 &W''(Y)&=2w_Y>0.
 \label{eq:wall-second-jet}
\end{align}
The strict inequality follows from the authenticated negative wall slope in
\cref{eq:consequence-profile-inputs}.  Analytically, a nontrivial solution
cannot have both $F_\star(x_{\rm w})=0$ and
$F_\star'(x_{\rm w})=0$, since uniqueness for the regular wall initial-value
problem would then give the vacuum solution.

For the present parameters,
\begin{equation}
 Y=\operatorname{arctanh}(1/5)=0.202732554054\ldots,
 \qquad N_{\rm w}=\frac{24}{25},
\end{equation}
and directed evaluation of the endpoint ledger gives
\begin{align}
 W''(Y)&\in[243.208686287,286.156494827],
 \notag\\
 A''(Y)&\in[1216.04343143,1430.78247414].
 \label{eq:certified-endpoint-ledger}
\end{align}
In particular, the leading wall amplitude is bounded away from zero.

\subsection{Boundary-controlled inertia-transform tail}
\label{sec:form-factor-tail}

In terms of \cref{eq:optical-weights}, define the finite optical transform of
the inertia density
\begin{equation}
 \cN(p)
 =\frac1p\int_0^Y A(y)\sin(py)\dd y
  -\int_0^Y W(y)\cos(py)\dd y.
 \label{eq:exact-inertia-transform}
\end{equation}
This is the profile-dependent numerator used in rotational matter-current
form-factor constructions.  We leave it unfiltered so that the result below
isolates exactly what follows from the certified inertia distribution.  This
continuum expression, rather than a finite radial quadrature, is used at large
frequency.

The relation between endpoint order and the asymptotics of finite Fourier or
Hankel transforms is standard; see, for example,
\cite{wang2018,okadayamane2024}.  We make no priority claim for that general
mechanism.  The result specific to this work is that the nonlinear profile
certificate authenticates the quadratic wall zero, the nonzero coefficient in
\cref{eq:certified-endpoint-ledger}, and the global consequence below.

\begin{proposition}[Sharp authenticated inertia-transform tail]
\label{prop:sharp-transform-tail}
For the exact certified profile,
\begin{align}
 \cN(p)
   &=W''(Y)\frac{\sin(pY)}{p^3}+o(p^{-3}),
 \notag\\
 \cN'(p)
   &=Y W''(Y)\frac{\cos(pY)}{p^3}+o(p^{-3}),
 \notag\\
 \cN''(p)
   &=-Y^2W''(Y)\frac{\sin(pY)}{p^3}+o(p^{-3}).
 \label{eq:sharp-p-minus-three-tail}
\end{align}
Moreover, the certificate gives
\begin{equation}
 (n_0,n_1,n_2)=
 (3.7575\!\times\!10^{10},
  3.7830\!\times\!10^{10},
  7.5660\!\times\!10^{10})
 \label{eq:inertia-transform-tail-constants}
\end{equation}
such that
\begin{equation}
 \abs{\cN^{(k)}(p)}\leq n_kp^{-3},
 \qquad p\geq1.
 \label{eq:global-inertia-transform-tail-bound}
\end{equation}
\end{proposition}

\begin{proof}[Proof sketch]
For $m=0,1,2$, set
\begin{equation}
 M_m^W=\norm{\partial_y^3(y^mW)}_{L^1(0,Y)},
 \qquad
 M_m^A=\norm{\partial_y^3(y^mA)}_{L^1(0,Y)}.
 \label{eq:third-derivative-norms}
\end{equation}
The exact interval-AD certificate, with a separate regular-origin Volterra--Lie
calculation, proves
\begin{align}
 (M_0^W,M_1^W,M_2^W)
 &\leq(2.5776861727020434\!\times\!10^8,
       6.94657005661323\!\times\!10^6,
       2.98269008103853\!\times\!10^5),
 \notag\\
 (M_0^A,M_1^A,M_2^A)
 &\leq(3.731716425863237\!\times\!10^{10},
       5.056441300063133\!\times\!10^8,
       1.391944736624204\!\times\!10^7).
 \label{eq:certified-third-derivative-norms}
\end{align}
These conservative values are upper displays of exact rational bounds.

Three integrations by parts are performed on the closed half interval
$[0,Y]$.  The endpoint jets in
\cref{eq:origin-optical-jets,eq:wall-optical-jets} cancel the lower-order
boundary terms but retain the wall term proportional to $W''(Y)$.  This point
is essential: the zero extension of $W$ has a jump in its second derivative at
$Y$ and must not be treated as a global $W^{3,1}$ function.  The six norms in
\cref{eq:certified-third-derivative-norms} bound the integral remainders and,
after differentiating under the integral sign, give
\cref{eq:global-inertia-transform-tail-bound}.  Keeping the wall terms and
applying the Riemann--Lebesgue lemma to the remainders gives the three sharp
expansions in \cref{eq:sharp-p-minus-three-tail}; derivatives retain the same
power because they differentiate the endpoint phase $pY$.

The same explicit bookkeeping gives the three outward-rounded constants in
\cref{eq:inertia-transform-tail-constants}.  In the combined archive they are
the upper endpoints of
\texttt{numerator\_derivative\_coefficients}; no form-factor prefactor or
frequency filter is applied to \(\cN\).
\end{proof}

The endpoint and six-norm certificate has SHA-256 digest
\begin{center}
\small\nolinkurl{1d5fe53786cc280006d7b1092d360556d4d8d8684e5ae3356ce8cd6d084e72a9}.
\end{center}
It binds the wall slope, inertia, endpoint ledger, all six derivative norms,
and the continuum tail envelope to one certificate identity.

\begin{remark}[Scope of the spectral statement]
The $p^{-3}$ coefficient is a property of this Dirichlet-confined inertia
profile.  A generic finite-pinning profile with nonzero wall value loses the quadratic
cancellation and can have a slower tail; no smooth-confinement universality
theorem is asserted.  The transform \eqref{eq:exact-inertia-transform} is defined
directly from the certified inertia density.  Identifying it with a detector
bath factor would require a specified current, interaction, switching
protocol, and open-system limit; none of those is asserted here.
\end{remark}
